\chapter{Proofs, finite certificates, and numerical evidence}
\label{app:verification}

A construction that uses exact algebra, analytic theorems, finite
computations, formalization, and numerical experiments needs more
than one notion of verification. This appendix explains those
notions without asking the reader to reconstruct the entire
research workspace.

\section{An exact finite certificate has a mathematical target}
\index{certificate!exact finite}

A finite lattice calculation can certify a vector equality
$a_{(n)}b=\sum_jc_jv_j$, a membership in a specified span, the
rank of a matrix, positivity of a Gram matrix, or the cyclicity
of a finite-dimensional top. The target must name the actual
vectors, mode conventions, coefficient field, and basis.
Otherwise the phrase ``the check passed'' leaves the mathematical
claim unclear.

The exact construction uses the lattice oscillator formulas and
algebraic coefficient fields such as $\Q(\sqrt{13},\sqrt{-3})$.
An equality in that field is stronger than a small floating
residual. A saved coordinate relation can be checked by a small
independent program without recomputing the original mode product.
That verifies the relation among the supplied vectors. It does
not independently verify that those vectors were correctly
computed as the claimed lattice products. Both checks are useful,
and their scopes should be distinguished.

The finite-closure theorem then converts an explicitly stated
family of finite identities into an all-weight strong-generation
conclusion. This implication is a mathematical proof, not an
extrapolation from a list of low-weight dimensions. In particular,
the provenance of recursively defined higher fields is part of
the hypotheses: a vector with the right weight is not enough if
it has not been shown to belong to the algebra generated by the
original fields.

\section{A ledger of implications for the charge-eight construction}
\label{sec:verification-construction-ledger}
\index{construction inputs!proof dependencies}

Inputs apply row by row; the diagonal net route branches before
the final fusion-positivity step.
Finite construction inputs retain the status explained above; the
structural implications do not independently replay them.
Write $H=V_{E_8}$ for the original ambient extension and $H'$ for
the later selected extension.

\begingroup
\small
\renewcommand{\arraystretch}{1.12}
\begin{longtable}{@{}>{\raggedright\arraybackslash}p{.15\textwidth}
 >{\raggedright\arraybackslash}p{.38\textwidth}
 >{\raggedright\arraybackslash}p{.41\textwidth}@{}}
\toprule
Step & Needed input & Output and logical boundary\\
\midrule
\endhead
Finite closure &
Specified PCT-stable lattice fields and ambient stress;
actual memberships and product identities;
recursive provenance of higher generators. &
Theorem~\ref{thm:c-recursive} gives all-weight strong generation.
The stated finite identities give properness and $\dim U_1=6$,
\eqref{eq:c-properness}. A few ranks alone do not.\\[5pt]

Positive stresses &
Actual PCT-stable embedded regular affine, Potts and parafermion
algebras; their verified projection identities. &
\eqref{eq:c-finite-positive} expresses $\omega_U$ as a positive
finite sum of their stresses. This is an energy identity,
not a commuting tensor decomposition.\\[5pt]

Strong rationality &
That stress family, compatible component tops, vacuum splitting,
and the finite-category rigidity argument. &
$C_2$-cofiniteness and vacuum projectivity
(Propositions~\ref{prop:c-c2}, \ref{prop:c-projective}) lead to
semisimplicity and strong rationality
(Proposition~\ref{prop:c-projective-unit};
Chapter~\ref{ch:regularity}). Positivity of every $U$-module
has not been assumed.\\[5pt]

Algebraic branching &
Strong rationality; finite core-top bounds, extension branching,
rank-four exclusion and indicator moments. &
$H\downarrow_U=\one\oplus X\oplus2Y$,
\eqref{eq:c-H-branching}; the boundary ring is Haagerup and there
are twelve chiral simples. This does not identify $Y=a$
or determine the unitary associator
(Chapter~\ref{ch:branching}).\\[5pt]

Complete unitarity &
Strong locality; independent algebraic and net sector counts;
common domains and paired-annulus comparison; fusion positivity
and unitary category classification. &
Positive modules and canonical fusion forms; then
$\Rep(U)\simeq Z(\HH_3)$ as unitary ribbon categories.
The original boundary category remains
$\HH_2$ or $\HH_3$
(Chapter~\ref{ch:unitarity}, especially
\eqref{eq:unitarity-pair-comparison}).\\[5pt]

Selected extension &
The center equivalence and the positive canonical Lagrangian
algebra belonging specifically to $\HH_3$; positive extension
and charge-eight holomorphic classification theorems. &
A selected extension
$H'\downarrow_U=\one\oplus b\oplus2a$, with $H'\cong V_{E_8}$.
That ambient isomorphism does not identify the two embeddings
of $U$ (Chapter~\ref{ch:full_theories}).\\[5pt]

Full local theories &
Diagonal: complete rationality of the net and Longo--Rehren.
Extension: the chosen positive algebra and extension/sewing theorems. &
The diagonal theory and $H'\otimes\overline{H'}$ exist, with
different Hilbert spaces, \eqref{eq:full-diagonal} and
\eqref{eq:full-E8}. Modular invariance alone would not establish
either sewing construction
(Chapter~\ref{ch:sewing}).\\
\bottomrule
\end{longtable}
\endgroup

Characters follow from the modular representation, admissible
principal parts and six-current coefficient, with uniqueness supplied by
Proposition~\ref{prop:char-unique}. They count states in the constructed
sectors; they do not replace the earlier domain or positivity
arguments. None of the rows identifies the spin-chain infrared limit.

\section{Reduction modulo a prime has a direction}
\index{finite-field rank test}

Suppose a matrix has entries in an algebraic number field and
all denominators remain invertible under a chosen reduction to
a finite field. If a minor is nonzero after reduction, it was
nonzero before reduction. Therefore
\[
 \rank_{\text{number field}}M\ge\rank_{\text{finite field}}\bar M.
\]
A nonzero modular minor proves a characteristic-zero rank lower
bound. A zero modular determinant does not prove a zero original
determinant: reduction can lower rank. Nor does a large modular
rank by itself prove that a target belongs to a span. A membership
needs an exact expression or an independent upper-bound argument
that closes the dimension count.

The choice of reduction matters. Algebraic generators must satisfy
their defining polynomials in the finite field, and denominators
must be checked. For example a map sending $\sqrt{13}$ to an
element whose square is not thirteen is not a valid specialization
of the intended coefficient ring, regardless of how plausible
the resulting rank looks.

\begin{example}[A modular dependence that disappears in characteristic zero]
\label{ex:verification-modular-rank}
Let $R=\Z[1/2,\sqrt{13}]\subset K=\Q(\sqrt{13})$ and take
\[
 A=\begin{pmatrix}1&\sqrt{13}\\\sqrt{13}&16\end{pmatrix},
 \qquad \det A=3.
\]
Specializing the coefficient ring to $\mathbb F_3$ by
$\sqrt{13}\mapsto1$ is valid: the defining square relation
is preserved and the denominator two is invertible. It sends
$A$ to the all-ones matrix, of rank one.
Its two columns coincide modulo three.
Over $K$, expressing the second column as a multiple of the
first would instead force the coefficient to be $\sqrt{13}$.
The lower coordinate would then be $13$, while the required
coordinate is $16$. Thus this modular span membership does
not lift to characteristic zero.

The specialization $R\to\mathbb F_{17}$ sending
$\sqrt{13}$ to $8$ is also valid, since $8^2=64\equiv13$.
Now the determinant remains $16-64\equiv3\ne0$, so the
specialized rank is two. The nonzero minor proves rank two
over $K$ as well. The prime three was admissible for the
coefficients but lost rank for this particular matrix.
These are maps from $R$, not field homomorphisms
$K\to\mathbb F_p$.

Finally, positive definiteness is a separate real statement.
The leading principal minors of $A$ are $1$ and $3$, so
Sylvester's criterion proves positivity under either real
embedding of $K$. Neither finite-field rank calculation
establishes that positivity.
\end{example}

\section{What a proof assistant theorem actually assumes}
\index{formal verification}\index{trust boundary}

A kernel-checked theorem can be highly informative while still
being conditional. Consider the abstract theorem
\[
 (H:P)\longmapsto H.
\]
It proves $P$ from a hypothesis $P$ using no unproved axiom in
its body. Inspecting only the theorem's global axiom list would
not tell us that an instance of $P$ is still needed.
Likewise, a formal dependency theorem can correctly compose
finite certificate assumptions, external results, and manuscript
lemmas without constructing every assumption for the intended
lattice VOA.

The repository's formalization contains real concrete mathematics:
mode identities, finite current and field relations, parts of the
properness argument, character and category arithmetic, and
logical composition of the construction. It also retains explicit
external, certificate, and analytic hypotheses. Some conclusion
nodes are concrete propositions about constructed objects;
others encode the intended theorem as a proposition parameter.
These distinctions are recorded in its trust-boundary document.

Accordingly, this textbook does not call the complete Haagerup
CFT construction machine verified. Formalizing a finite identity,
formalizing the implication from that identity, and supplying a
complete formal instance of the whole theory are different tasks.
The available formal results help locate the remaining work
precisely.

There is also a separate, stronger recorded formal result about
the uniform coefficient construction in
Chapter~\ref{ch:arithmetic}. The theorem
\sourcefile{HaagerupIzumi.exists_EGEquations_odd}
constructs complex coefficient data satisfying the complete
Evans--Gannon equations for every odd natural number. The
underlying hyperbolic gamma interface is constructed from its
regularized integral; the final theorem does not assume an
unproved special-function interface. The recorded build and
transitive axiom audit in \sourcefile{lean/README.md} report
only Lean's standard logical axioms. The corresponding coefficient
source revision is
\sourcefile{698b66ac2c667e68b1e19c196fb91e84a0c82235}.

This all-order coefficient theorem is not a finite numerical
spot check. It is also not a formal theorem that reconstructs a
fusion category, proves a positive adjoint structure, or constructs
a CFT. Those later steps have separate mathematical inputs.
For this book the final declarations, constructors, source
status, and recorded audit were inspected without rebuilding
the Lean development. The short revised analytic exposition
has documented domain gaps that the expanded/formal development
addresses; one should use the latter when following the complete
analytic proof. Keeping this result separate from the conditional
$c=8$ formalization avoids both understating the coefficient proof
and overstating its categorical conclusions.

\section{Numerical precision and continuum control}
\index{residual bound}\index{continuum limit!evidentiary scope}

For a normalized approximate eigenvector $v$ of a Hermitian
finite matrix $H$, a residual
$r=\norm{Hv-\lambda v}$ implies that some exact eigenvalue lies
within $r$ of $\lambda$. It does not identify which ordered
eigenvalue it is without further isolation information, and it
does not give a finite-size error bound relative to a continuum
dimension.

Similarly, a variational MPS energy is an upper bound within the
physical variational setting. A small gradient establishes
stationarity in the stated coordinate chart and tolerance; it
does not establish a global minimum. A transfer eigenvalue close
to the unit circle can arise from critical correlations or from
multiple symmetry-related components. A high overlap of a
source with one vector controls that source action on the vacuum,
not its matrix elements on every excited state.

The chain chapters retain the completed finite results and state
their continuum assumptions locally. They do not convert decimal
precision into asymptotic certainty. The open infrared problem
requires both a candidate full theory and controlled comparisons
of dimensions, symmetry sectors, and normalized local data.

\section{Reproducing examples responsibly}

Most examples in this book need only hand algebra, short exact
polynomial manipulations, or existing finite tables. Reproducing
them does not require rerunning the largest lattice calculations.
The accompanying \sourcefile{check_examples.py} reproduces
$S^2=I$ and all $1,728$ Verlinde coefficients of the displayed
Haagerup modular matrix in exact cyclotomic arithmetic. It also
checks two twisted-character vectors and selected descendant,
Gram-matrix, and correlator examples. Its recorded successful run
has this finite scope; it does not certify the lattice construction.
The separate \sourcefile{check_order_three.py} verifies the nine
quadratics and all $81$ cubics of
Example~\ref{ex:HI-n3-cubics} using four-integer normal forms.
The hand reduction in that example explains why the cancellations
occur, independently of the program's finite enumeration.
The third checker, \sourcefile{check_core_tops.py}, reconstructs
all fourteen core tops in Chapter~\ref{ch:branching}. It keeps every
tied charge and its multiplicity, checks the finite search cutoff,
and verifies the trace certificate for the $17/10$ bound coefficient
by coefficient. This confirms the displayed arithmetic under the
core-realization and compatible-top hypotheses; it does not establish
those representation-theoretic hypotheses by running a program.
The Ising calibration figure uses only the closed formulas proved
in its worked example, with no numerical diagonalization or fit.

Before any substantial numerical calculation, estimate all
simultaneously live allocations: basis vectors, sparse assembly,
eigensolver workspaces, temporary copies, and tensor contractions.
The size of the final result is not a bound on peak memory.

For the shared research workspace accompanying this text, the
canonical machine policy and central research launcher govern
numerical admission. Historical scripts and resource exceptions
are provenance, not current authorization. The textbook's PDF
builder is a bounded typesetting tool with shell escape disabled;
it does not execute scientific calculations or replay certificates.

The source files for the book are separated from the research
sources. Document compilation checks references, layout, and
syntax. A mathematical review checks equations and implications.
Neither activity silently upgrades an unverified scientific input.
Keeping those responsibilities explicit makes the book a reliable
starting point for both learning and further research.
